\documentclass[tikz,border=8pt]{standalone}
\usepackage{fontspec}
\usepackage{xeCJK}
\IfFontExistsTF{Segoe UI}{\setmainfont{Segoe UI}}{\setmainfont{Arial}}
\IfFontExistsTF{Microsoft JhengHei}{\setCJKmainfont{Microsoft JhengHei}}{\IfFontExistsTF{Noto Sans CJK TC}{\setCJKmainfont{Noto Sans CJK TC}}{\setCJKmainfont{Noto Sans TC}}}
\IfFontExistsTF{Microsoft JhengHei}{\setCJKsansfont{Microsoft JhengHei}}{\IfFontExistsTF{Noto Sans CJK TC}{\setCJKsansfont{Noto Sans CJK TC}}{\setCJKsansfont{Noto Sans TC}}}
\usetikzlibrary{arrows.meta,positioning,calc}
\definecolor{paper}{HTML}{FFFDF8}
\definecolor{navy}{HTML}{0F172A}
\definecolor{muted}{HTML}{334155}
\definecolor{blue}{HTML}{2563EB}
\definecolor{bluesoft}{HTML}{DBEAFE}
\definecolor{gold}{HTML}{D97706}
\definecolor{goldsoft}{HTML}{FEF3C7}
\definecolor{red}{HTML}{DC2626}
\definecolor{redsoft}{HTML}{FEE2E2}
\definecolor{line}{HTML}{CBD5E1}
\begin{document}
\begin{tikzpicture}[x=1cm,y=1cm,>=Latex]
\path[fill=paper] (0,0) rectangle (16,9.6);
\node[font=\bfseries\fontsize{18}{22}\selectfont,text=navy] at (8,8.82) {預算時間封鎖線};
\node[font=\fontsize{10.5}{13}\selectfont,text=muted] at (8,8.12) {做預測時，只能使用當天真的拿得到的資料};

\draw[rounded corners=10pt,fill=bluesoft,draw=blue,line width=1pt] (.95,2.55) rectangle (6.85,7.1);
\node[font=\bfseries\fontsize{12}{15}\selectfont,text=navy] at (3.9,6.55) {預算當天可用};
\node[font=\fontsize{8.5}{11}\selectfont,text=muted] at (3.9,6.17) {可以放進模型的資料};
\foreach \y/\t in {5.45/歷史銷售,4.82/價格與促銷計畫,4.19/通路庫存,3.56/已知季節因素}{
  \draw[rounded corners=5pt,fill=white,draw=line,line width=.9pt] (1.65,\y-.25) rectangle (6.15,\y+.25);
  \node[font=\fontsize{9.2}{11}\selectfont,text=navy] at (3.9,\y) {\t};
}
\foreach \y in {5.45,4.82,4.19,3.56}{\fill[blue!75!navy] (1.38,\y) circle (.07);}

\draw[dashed,line width=1.25pt,draw=navy] (8,2.25)--(8,7.35);
\draw[rounded corners=12pt,fill=navy,draw=navy] (7.35,7.15) rectangle (8.65,7.58);
\node[font=\bfseries\fontsize{9}{11}\selectfont,text=white] at (8,7.37) {預算日};

\draw[rounded corners=10pt,fill=redsoft,draw=red,line width=1pt] (9.15,2.55) rectangle (15.05,7.1);
\node[font=\bfseries\fontsize{12}{15}\selectfont,text=navy] at (12.1,6.55) {事後才知道};
\node[font=\fontsize{8.5}{11}\selectfont,text=muted] at (12.1,6.17) {不能偷渡進訓練資料};
\foreach \y/\t in {5.45/實際銷售結果,4.82/退貨與缺貨,4.19/競品突然降價,3.56/月末庫存異常}{
  \draw[rounded corners=5pt,fill=white,draw=line,line width=.9pt] (9.85,\y-.25) rectangle (14.35,\y+.25);
  \node[font=\fontsize{9.2}{11}\selectfont,text=navy] at (12.1,\y) {\t};
}
\foreach \y in {5.45,4.82,4.19,3.56}{
  \draw[red,line width=1.2pt] (9.58,\y-.08)--(9.74,\y+.08);
  \draw[red,line width=1.2pt] (9.74,\y-.08)--(9.58,\y+.08);
}

\draw[draw=line,line width=.9pt] (.95,1.55)--(15.05,1.55);
\node[font=\fontsize{10.5}{13}\selectfont,text=navy] at (8,1.02) {先把時間線封好，模型準不準才有討論的意義};
\end{tikzpicture}
\end{document}


